\Needspace{12\baselineskip}
\subsection{Volume tests and threshold tests}
\label{sec:radial-gauduchon-test-limits}

The volume quotient divides a linear pairing by a determinant integral; the bigness threshold divides it by a trace. The same metrics can approach both infima, but need not do so. The radial model distinguishes these possibilities and gives actual smooth balanced tests concentrating on the degenerate fiber.

Keep the data of \cref{thm:radial-positive-volume-nonbig}, and put
\[
 T_\epsilon=\gamma_\epsilon^m,\qquad
 I_\epsilon=\int_B(a+\epsilon)^{-2}\omega^m,
 \qquad M_\epsilon=m+I_\epsilon/H_\epsilon,
 \qquad g_\epsilon=M_\epsilon^{-1/m}\gamma_\epsilon.
\]
Write $Z=\{p\}\times\mathbb P^1$. Its integration current has $\int_Zh=1$ and annihilates $\theta$.

\Needspace{18\baselineskip}
\begin{proposition}\label{prop:radial-volume-test-limits}\index{Gauduchon variational volume!normalized radial tests}
The metrics $g_\epsilon$ are balanced, have $\int_Xh\wedge g_\epsilon^m=1$, and satisfy
\begin{equation}\label{eq:radial-test-pairings}
 \int_X\theta\wedge g_\epsilon^m
   =\frac{n-\epsilon I_\epsilon/H_\epsilon}{M_\epsilon},
 \qquad D_\mu(g_\epsilon)
   =\frac{(n/H_\epsilon)^{m/n}H_\epsilon}{M_\epsilon}.
\end{equation}
Their volume quotients tend to $n^{1/n}$. If $n=3$, then $g_\epsilon^2$ converges weakly to $[Z]$, and
\begin{equation}\label{eq:radial-threefold-test-collapse}
 \int_X\theta\wedge g_\epsilon^2
       \sim\frac6{\log(1/\epsilon)},
 \qquad D_\mu(g_\epsilon)
       \sim\frac{2\,3^{2/3}}{\log(1/\epsilon)}.
\end{equation}
If $n\geq4$, put $M_0=m+(m-1)^2/[m(m-2)]$. Then
\[
 \int_X\theta\wedge g_\epsilon^m\longrightarrow\frac n{M_0}>0,
 \qquad D_\mu(g_\epsilon)\longrightarrow\frac{n^{m/n}}{M_0}>0.
\]
Thus this family approaches the zero bigness threshold precisely in its three-dimensional case.
\end{proposition}
\begin{proof}
The power formula in the preceding theorem gives
\[
 T_\epsilon=\frac{\omega^m}{H_\epsilon(a+\epsilon)^2}
               +mA_\epsilon^{m-1}\wedge w.
\]
Closed Stokes and $[A_\epsilon]=[\omega]$ give
$\int_B\omega\wedge A_\epsilon^{m-1}=1$. Consequently the two terms have $h$-traces $I_\epsilon/H_\epsilon$ and $m$. This proves the trace normalization. Subtracting $\epsilon w$ from $\theta_\epsilon$ subtracts $\epsilon I_\epsilon/H_\epsilon$ from its pairing $n$. Rescaling by $M_\epsilon^{-1/m}$ divides both the linear pairing and $D_\mu$ by $M_\epsilon$, proving \eqref{eq:radial-test-pairings}.

Polar integration gives
\[
 I_\epsilon=m t^2\int_0^1
              \frac{x^{m-1}}{(x+t\epsilon)^2}\,\mathrm{d}x.
\]
When $m=2$, put $q=\epsilon/2$. Direct integration yields
\[
 I_\epsilon=\tfrac12\left(\log\frac{1+q}{q}
                      +\frac q{1+q}-1\right),
 \qquad H_\epsilon=1-q\log\frac{1+q}{q}.
\]
Hence $I_\epsilon\sim\tfrac12\log(1/\epsilon)$ and $\epsilon I_\epsilon\to0$. Since $H_\epsilon\to1$, the pairings in \eqref{eq:radial-threefold-test-collapse} follow.

The mixed positive term $2A_\epsilon\wedge w$ has $h$-trace two, so its quotient by $M_\epsilon\to\infty$ tends weakly to zero. The horizontal probability
\[
 I_\epsilon^{-1}(a+\epsilon)^{-2}\omega^2
\]
concentrates at $p$: outside any fixed neighborhood its unnormalized density is bounded while $I_\epsilon\to\infty$. Its mass is one. For a smooth $(1,1)$-test form on $X$, integration along the projective-line fibers gives a smooth function on $B$. The probability convergence therefore gives the horizontal current limit $[Z]$, with coefficient $I_\epsilon/(H_\epsilon M_\epsilon)\to1$.

For $m\geq3$, monotone convergence instead gives
\[
 I_\epsilon\longrightarrow m t^2\int_0^1x^{m-3}\,\mathrm{d}x
                =\frac{(m-1)^2}{m(m-2)}.
\]
In both cases $\epsilon I_\epsilon\to0$. Using $H_\epsilon\to1$ in \eqref{eq:radial-test-pairings} proves the remaining assertions and the volume-quotient limit.
\end{proof}

The higher-dimensional family still minimizes the perturbed volume and approaches the endpoint volume. Its positive trace-normalized pairing does not contradict nonbigness: a different family approaches the bigness threshold.

\Needspace{16\baselineskip}
\begin{proposition}\label{prop:radial-localized-balanced-tests}\index{Gauduchon variational volume!localized balanced tests}
For every $n\geq3$ there are smooth balanced metrics $b_\rho$, $0<\rho<1$, with
\[
 \int_Xh\wedge b_\rho^m=1,\qquad
 b_\rho^m\rightharpoonup[Z],\qquad
 \int_X\theta\wedge b_\rho^m\longrightarrow0,
 \qquad D_\mu(b_\rho)\longrightarrow0
\]
as $\rho\downarrow0$.
\end{proposition}
\begin{proof}
Choose a nonzero smooth $f_\rho\geq0$ on $B$, supported where $x<\rho$ and normalized by $\int_Bf_\rho\omega^m=1$. The form
\begin{equation}\label{eq:radial-localized-test-form}
 \Gamma_\rho=\frac{f_\rho\omega^m+\rho h^m}{1+n\rho}
\end{equation}
is smooth strictly positive and closed: the first summand has base top degree, and $h$ is K\"ahler. The smooth positive root argument accompanying \cref{thm:lamari} gives a unique metric $b_\rho$ with $b_\rho^m=\Gamma_\rho$. Closedness makes it balanced. Since $\int_Xh^{m+1}=n$, its $h$-trace is one.

Put $C_\theta=\int_X\theta\wedge h^m$. Then
\[
 0\leq\int_X\theta\wedge\Gamma_\rho
 \leq\frac{\rho/t+\rho C_\theta}{1+n\rho}\longrightarrow0.
\]
The lower bound of \cref{thm:radial-positive-volume-nonbig} gives
$D_\mu(b_\rho)\leq n^{-1/n}\int_X\theta\wedge\Gamma_\rho$, proving determinant collapse. The probabilities $f_\rho\omega^m$ concentrate at $p$ because their supports shrink to that point. Fiber integration as in the preceding proof gives the current limit $[Z]$; the term $\rho h^m$ tends to zero.
\end{proof}

Thus $[Z]$ belongs to the weak closure of actual smooth balanced normalized tests. It annihilates $\theta$ and has zero determinant-weighted measure by \cref{prop:nef-annihilator-determinant}. These two explicit families realize the separation in \cref{prop:gaud-positive-volume-collapse} without a general smoothing assertion for Gauduchon currents.
